









\section{old discussion section}




% Future research considering multiple linguistic features together could be especially illuminating here; it is possible that the Alpaca and Squirrel do also make changes towards more hegemonically feminine indexicalities with the cisgender listener, but that these changes are along features other than /s/, such as overall pitch (f0), intonation, voice quality, or discourse markers.

% It is notable that these attitudes and concerns seemed  even though all four speaker participants did express concerns about cisgender judgement and that they would change how they speak to be more feminine around cis listeners -- evidencing mental salience of the cisgender listening subject --

% only two out of the four speakers implement this change in practice with /s/. Future research considering multiple linguistic features together could be especially illuminating here; it is possible that the Alpaca and Squirrel do also make changes towards more hegemonically feminine indexicalities with the cisgender listener, but that these changes are along features other than /s/, such as overall pitch (f0), intonation, voice quality, or discourse markers. I return to the issue of investigating only one linguistic variable below.

% In the model I initially constructed, speakers who expressed greater degrees of struggling with voice dysphoria, greater desire have a voice which 'passes' as the voice of a cisgender woman, or a greater sense of danger and pressure to conform around cisgender listeners would be more likely to shift to a fronter /s/ pronunciation. On the other hand, speakers who expressed stronger attitudes of resistance against cisnormativity, differentiated themselves from hegemonic femininity, or were concerned that their voice might make them sound like a gay man would be more likely not to shift to a fronter /s/ pronunciation. In interpreting the explanatory failure of this model on my data, it is first important to note that its formulation tacitly constructs a mutual exclusivity of the former and latter concerns, which is not upheld by my data. It seems that for Mantis and Jellyfish, the former concerns (which I thought would predict a shift to fronter /s/) are caused by the latter concerns (which I thought would predict no shift) -- it was discomfort with hegemonic feminine norms and concerns with gay-male misrecognition which caused these speakers to feel pressure to conform around cisgender listeners. In the discussion below, I return to the literature to reconsider the predictive factors I initially proposed in light of the patterns which actually emerged, specifically focusing on interpreting my results through the lens of indexicality.

% The issue of misrecognition is central here. In Mantis' discussion of how her voice is percieved, sounding like a straight man is not considered a possibility; if she is perceived correctly she will sound like a queer woman, and if she is perceived incorrectly she will sound like a gay man. Within the model of /s/ indexicality brought forth by \textcite{calder_sissy_2019} (recall figure \ref{figure_fractal_recursivity_indexical_orders}), other semiotic information is used to determine through which indexical order /s/ pronunciations will be interpreted. From this perspective, a shift in /s/ pronunciations could be considered useless for Mantis, in that whether she will be heard as a gay man or a cis woman cannot be determined by her /s/ pronunciation. Yet she shifts, contradicting the model. The extent of her shift provides help here; she shifts to /s/ pronunciations which fall in closer line with those expected of cis women than those reported for cis gay men. Within the literature reviewed, \textcite{zimman_hegemonic_2013} gives the highest mean CoG for cisgender gay men (not in drag) in the literature reviewed at ~6550Hz, whereas Mantis' /s/ pronunciations with the Cis CP had a mean CoG of 7964Hz, well within the hegemonically expected range of cis women. In \textcite{calder_sissy_2019}, SoMa queens front /s/ further than Castro queens in order to forge a higher indexical order and index a radical queer femininity beyond mainstream cisgender gay male femininity. Here, I argue that Mantis fronts /s/ further than cis gay men do, but that no higher order of indexicality can be forged, as she must make do with the two hegemonic indexical orders already existing in the assumed indexical inventory of her listener; instead, she uses this much more fronted hegemonically-feminine /s/ access to the lower level of indexicality and encourages her /s/ pronunciations to be interpreted as those of a woman on the $n$th level, and not those of a gay man on the $n + 1st$.

% Although this extension of the model works for Mantis, the situation is different for Jellyfish. If it is indeed the case, as Jellyfish mentions, that ``when you're a lesbian sounding like a gay guy and like having a faggcent will get you clocked way more than sounding like boy,'' why does she choose to shift to a fronter /s/ rather than not shifting /s/ and 'sounding like a boy'? Although it is not easy to answer this question with certainty from the data available and without context of other indexical features, I see three possible explanations. First, it is possible that Jellyfish is simply mirroring the /s/ pronunciations of her CPs, and that questions of perception of gender and sexuality are less relevant to this case. A point in favor of this explanation is the priority Jellyfish lends to her lesbian identity; because the Cis CP was also a lesbian, it is possible that the function of linguistic mirroring in this case could have been a move to create greater similarity between herself and her CP, a move of `adequation' from the perspective of the `tactics of intersubjectivity` theorized by \textcite{irvine_language_2000}. The second possible explanation has to do with the complexity of Jellyfish's relationship with expressing boyness. Although she indicates that it is at least hypothetically safer for her to sound like a boy than a gay man around cis women, she also clarifies that ``I'm like less concerned about like trying to like overly feminize my voice when I'm around other trans people'', and that ``I love talking like a boy when I'm interacting with other boys''. From this perspective, it is possible that she felt less comfortable to sound like a boy with the Cis CP, and thus ``with like, it's mostly people I don't know, I'll like try harder to um have like a more conventionally feminine voice.'' A point of dissatisfaction with this second explanation is that it breaks down when considered for the other speakers; as mentioned above, all speakers indicated similar sentiments that they would shift towards hegemonically feminine speech with cisgender listeners, but Alpaca and Squirrel do not in fact shift /s/ in this way. In her discussion of her concerns with misrecognition, she underlines that although she is proud of ``being able to more or less manipulate [her voice] how [she] want[s]'', she nonetheless is sometimes bothered that she might
% \begin{quote}
% ``default back to talking like a gay guy around cis women ... It's kind of rolling the dice what voice comes out on a certain day, and sometimes the faggcent comes back''
%\end{quote}
%Her notion that a voice more prone to misrecognition (`the faggcent') is a `default' for her around cis women and her metaphor of styleshifting to rolling dice make it clear that she does not feel fully in control of her voice. Here, it is possible to understand her shift to a fronter /s/ pronunciation with the Cis CP as an example of this sort of `defaulting', rather than an agentive choice aligned with her intentions and goals. It is further interesting that in her discussion of how she talks around other trans people, comfort in her voice is also described in terms of defaulting:
%\begin{quote}
%``at a certain point, I like stop performing and sort of default back to like my neutral state which is still just like talking like a boy, which is fine, like I don't really feel much pressure to change that.''
%\end{quote}
%Although Jellyfish does acknowledge that ``it's all performance'', in line with feminist understandings of gender performativity, it is nonetheless notable that Jellyfish perhaps considers hegemonic femininity somewhat of a difficult or unnatural performance for her, and that it does not feel like her neutral state. Here, it is clear that styleshifting is not effortless, always available, and without cost; instead, it comes to be limited.

%From this, I am motivated to take up the question of the interplay between agency and limits on agency as central to the interpretation of my data. In the introduction to this thesis, I mentioned that agentic models of linguistic variation, which see speakers' linguistic practices as strategically and intentionally deployed to achieve social goals, was limited by the restrictions imposed by the importance of the listener's role in interpretation. From my results, it seems additionally important to consider the degree to which variants are available and natural to speakers; even if a speaker's vocal tract holds the potential to produce a wide degree of linguistic variants, it is nonetheless possible that a speaker may only be familiar with using only a smaller range of variants, or experience some variants as more effortful and less natural due to their varying familiarity. An analogy can be made to phonology in language learning; it is an uncontroversial fact of language that all humans with typical vocal tracts are capable of all phonetic and phonemic distinctions in all human languages, but it is also uncontroversial that learning the phonology of a new language is challenging and time-consuming, and that effort is required to sustain accurate pronunciation.

% I argue that this framework of limited agency helps explain Alpaca and Squirrel's lack of normative shifting as well. For Squirrel, the reasoning is fairly simple; as a transfem early in the process of transition, Squirrel considers a feminized voice as only accessible in the future. The inaccessibility of her desired voice -- her voice not being ``they way [she] would want it to be if [she] could like, pick it'' -- is complete enough that she does not use her voice at all on the internet, the sphere in which she wishes to be perceived as feminine. Being not fully out in in-person interactions means that she is for the time being continuing to assume she might be perceived as masculine, and does not yet feel able to combat these perceptions through the voice.

% For Alpaca, I believe that this framework can help explain her lack of a frontwards shift, although the backwards directionality of her shift remains difficult to explain. Alpaca mentions learning ``the exercise to raise your larynx'', and has the sense that she has implemented this into her everyday speech over the course of her transition, noting that she ``learned how to do it and then [she] found [her]self sort of subconsciously doing it'' over time. The status of agency in this example is instructive to unpack; despite the ability to raise the larynx always existing in her vocal tract, learning the exercise she mentions provides a greater degree of access to this variant. That she found herself ``subconsciously doing it'' indicates another complex relationship with agency, whereby changes could not be forced but become more natural and less effortful over time. With respect to /s/, it is possible that this resource for gendered indexicality is not easily accessible for Alpaca, and perhaps feels less needed in the context of her perceived success in feminizing her voice over the years of her transition. Because shifting /s/ fronter would require accessibility and effort, this change is not made. That this cannot account for why Alpaca shifts to backer /s/ pronunciations is not ideal, but perhaps not a problem for the framework; there do not seem to be any factors found in Alpaca's identity or interview which could be used to explain the directionality of her shift, so it is likely that more data would be needed to resolve this issue.

% It is also important to clarify that this framework of limitations on agency certainly does not imply that all of my speakers' behavior can only be explained by limitations rather than agentivity. As discussed above, I argue that Mantis represents an example of agency as a better explanation than limitation, in that her shift is the most drastic, and it simultaneously achieves both of her goals; to protect herself from dysphoria and negative judgment by her listener by taking up fronted /s/ as a hegemonically feminine index, and to protect herself from gay-male misrecognition by using fronter /s/ pronunciations than would be expected from gay men. My first possible explanation of Jellyfish's behavior may also suggest agency as a primary explanatory frame. These sorts of goal-satisfying use of linguistic resources is unproblematically read as agentive, even though it is defensive in Mantis' case. My argument here is for the both-and of agentivity and limitation, seeing that my data is best explained when both considerations are given weight and relevance. After all, this model more suitably describes the social world my participants live within; they are limited by their societies and their abilites, but nonetheless can agentically and strategically make choices when those choices are accessible and intentionally available.

% affective mediating layers / talk about ``customer service voice'' as ``not even an intentional gender thing''?

% Under this interpretation, it could be claimed that Mantis and Jellyfish are more practiced, so to speak, with speech forms with feminine indexicalities (perhaps both within and outside of the hegemonic indexical field) than Squirrel and Alpaca, even though the former two have been out as transgender and known they were trans for similar amounts of time as Alpaca. In this case, it is possible that shifting is more easily accessible to Mantis and Jellyfish than the former two speakers, as they are more practiced within the feminine range of /s/ pronunciations. This interpretation is possible,



% however, at that time, because they had otherwise male identifications, these features come to be interpreted as indexing gay male

% Returning to the possibility mentioned above that Mantis accomodates more to the

%Instead, I argue that social explanations motivations must be considered.
% it is necessary to evaluate whether my data can be explained through accomodation, without reference to the factor of trans status of the listener, the dynamic of the cisgender listening subject's hegemonic indexicality of /s/, or the trans background and intentions of the speakers. Is it possible to explain the differences in my speaker's behavior -- that Mantis and Jellyfish shift fronter, Alpaca shifts backer, and Squirrel does not shift -- in this way?
% It is at this point that CP's trans status emerges as a necessary explanatory factor;
% wouldn't it predict that transfems would accomodate more to other transfems than cisfems? Second, still is not able to differentiate between squirrel+alpaca and mantis+jellyfish.
% Mantis and Jellyfish's shifts are read as accomodation, why did they choose to accomodate where Squirrel and Alpaca did not? If it were simply the case that all speakers accomodate, these speakers would be expected to accomodate too.
% I nonetheless argue that accomodation should be considered an outcome, rather than a complete explanation. In other words, it is insufficient to state that speakers are accomodating without interpreting the motivations behind their accomodation.
% I argue that both of these frameworks must be considered as indexical moves rather than interpretive explanations. In the case of the first pattern, I problematize as a generalization which can only become an interpretation when normativity is assumed to be an explanation rather than an object of criticism-- it is not sufficient to say `she accomodates because she is a woman', without questioning what populations the generalization is based upon, and why women might accomodate more than men. Here, a reinterpretation through indexicality is needed. That accomodation research within cisgender heterosexual white populations has found women to accomodate more than men can be interpreted as an indexical relationship of co-occurrence; it is possible that the act of accomodating itself is indexical of femininity within the hegemonic indexical field.
% At this juncture, I argue that interpretation of my speakers' behavior must return to the social-indexical framework I set out in the literature review, confirming the relevance of trans status of the listener, the indexical structure of norms surrounding /s/ and accomodation, as well as the background and intentions of the speakers as transfems. An approach which does not incorporate these factors as relevant, simply relying on models of pure accomodation, cannot differentiate between these speakers meaningfully.
% In light of these considerations, I return to the question posed above: why Jellyfish and Mantis front /s/ with the Cis CP, accomodating and moving towards the hegemonic norm, and why Squirrel and Alpaca do not. Is it possible that Jellyfish and Mantis are just accomodators, and that the cisness of the listner doesnt matter? Well, I do not have any reason to believe that they intend to index femininity or positive stances towards their listeners more than Squirrel and Alpaca, so I don't think so. Also, they don't accomodate to the same extent to the trans listener as they do to the cis listener: Mantis's competition case.


% I argue she does not because she is not used to presenting feminine in face-to-face (not online) interactions, and has not yet learned (consciously or unconsciously) how to pronounce /s/ differently than she did pre-transition. \textcite{zimman_gender_2021} demonstrates that a multi-dimensional model of gender is needed to explain the /s/ pronunciation patterns of his gender-diverse set of speakers, including gender assigned at birth, current gender identity, and gender role -- what others know about a speaker's gender. My data agrees with Zimman's affirming the importance of considering gender role, since gender role likely provides an explanation for Squirrel's lack of a shift.





\section{other things}
``I am a woman but I'm not even that feminine ... my expression has gone from femme to butch, or futch a little bit". The invocation here of the lesbian in-group gender spectrum from 'femme' through 'futch' to 'butch' leads to the possibility of interpreting her mean [s] pitch as in line with cisgender queer womens', as would be backed up by the similarity of her mean [s] CoG to that of the cisgender queer women \textcite{hazenberg_language_2012} studies in Ottowa.

Although these observations may help to understand why Mantis' overall mean [s] CoG was somewhat higher than some of her peers, they are not particularly helpful to understanding her shifting behavior. Aside from these more confident and positive assertions of Mantis' identity, she also identified two areas of difficulty and concern when it came to


 It is possible that this high [s] is part of what has brought ``so many people'' to



“I've had so many people … tell me that they thought I was a cis woman with my voice”
“There would be gay men who have gay voice ... and I'd be like ‘oh I can't distinguish my voice from theirs’”
“especially with cishet women I feel … pressure for sure … it’s almost like a competition”

	
	

Queer people, in existing, make a revolutionary claim: sex, gender, and sexual orientation are not natural biological phenomena, and they do not always fall into one of two possible arrangements. That people exist outside of these cagetories and their expected binary arrangements -- male-thus-man-thus-desiring-women or female-thus-woman-thus-desiring-men -- exposes these categories and their order as social fictions. For example, the notion of a pansexual trans woman, whose body may be categorized `male', whose social gender is nonetheless `woman', and whose sexual desire is not limited to members of a single gender category, is troubling to this order.

\textcite{silverstein_langauge_2023} argues that categories of ientity, such as gender, arise through discourse. Here, discourse is taken broadly, from the smallest scale of everyday interaction up to the largest scales of cultural, political, and literary histories.



Despite the importance of this turn towards an agentic model, research which centers speakers' individual agency nonetheless runs the risk of erasing situations where individuals agency is limited and self-determination is denied or questioned. [the listening subject].


A complete model both highlights the agentivity of the speaker and acknowledges the injustices they face which may limit that agentivity; furthermore, a model which places the `doing' before the `doer' as in \textcite{butler_gender_1990} [PAGENUMBER], may be able to deconstruct self-other dialectic entirely and understand how it comes to be from its performative parts.

 nonetheless has the potential to erase the realities of the voice for many speakers; 


This arises  a history of sociolinguistic research which instead considered a speaker's behavior to be a product of their identity (\cite{eckert_three_2012}). This turn to acknowledge the individual's role in the construction of their identity is very important, especially with regard to transgender people. Trans people often consider gender to be a matter of self-identification and selfdetermination (\cite{zimman_linguistic_2017}), rather than a consequence of bodily characteristics; this is consonant with a view of the voice as a resource speakers agentically employ to construct their identities. Despite the importance of this agency-first approach as a move away from essentialist understandings, 




Sociolinguistics




\section{The Cisgender Listener}
The above claim that trans speech is often more scrutinized by cisgender people for gender conformity, that trans people are aware of this, and that their knowledge may manifest in linguistic practice is based upon existing findings in the field of transgender linguistics. Investigating transgender knowledge, \textcite{heffernan_constructing_2023} finds that transfeminine speakers characterize a `cisgender listening subject' in their metalinguistic commentaries which casts negative judgements when linguistic gender norms are not followed, and that transfeminine speakers expressed varied ways of negotiating this, from accomodating these expectations by a change in linguistic behavior towards conformance to asserting gender to be self-determined, resisting these judgements, and talking as one would prefer.  \textcite{jessee_cisgender_2025} emperically confirm that transgender speech is indeed subject to this greater gender scrutiny than cisgender speech, finding that cisgender listeners apply stricter gendered standards to speech when they are told the speaker is a transwoman than when told the speaker is a ciswoman. Although \textcite{parnell-mooney_transgender_2019} predates these more recent findings, he presents a design very similar to my own but with a focus on transmasculine and nonbinary speakers, in which it was found that these speakers do in fact pronounce /s/ with different pitch when talking to a trans person versus a cis person, but that variation was did not pattern across speakers.  

\section{Transfeminine Speakers in Sociophonetics}
My choice to focus on transfeminine speakers is motivated by three reasons: the relative scarcity of sociolinguistic research which investigates transfeminine speech on the phonetic level, the emperical gap left by left by the exclusion of transfeminine speakers from \textcite{parnell-mooney_transgender_2019}, and the contributions to sociolinguistic theory, especially existing questions in the scholarship on /s/, that may arise by investigating transfeminine speech.  The scarcity of sociophonetic research on transfeminine speakers has been noted by \textcite{zimman_transgender_2020}; in a review of the place of trans people (and other groups which fall outside of the gender binary) in the language and sex, gender, and sexuality literature, they find that  in the Global South, research has primarily focused on third and fourth-gender groups  which bear similarities to transfeminine people, and in the Global north it has leaned towards transmasculine people, ``leaving trans feminine people in the Global North and trans masculine people in the Global South underrepresented.'' As I turn to the reasoning behind my choice to look at /s/, I will demonstrate this gap and discuss the ways in which a study of transfeminine interlocutor-oriented style-shifting (the sociolinguistics term for changes in an individual's speech depending on who they're talking with, see eg. \cite{rickford_addressee_1994}) could contribute new insights on existing theoretical questions raised by prior work on /s/.

\section{/s/ in Language and Gender Research}

Termed `sexuality's most iconic variable' by \textcite{calder_gay_2020}, /s/ has become one of the most well-studied sites of socially-relevant variation in speech (called `variables' in sociolinguistics) in the field of language and sex, gender, and sexuality. This scholarship has primarily oriented itself around finding differences in /s/ pronunciation and perception between sex, gender, and sexuality groups, individuals, and contexts, and attempting to explain those differences. The most commonly studied aspect of /s/ is frontness, an analouge of /s/ pitch; a fronter tongue position produces a higher-pitch (see \cite{fuchs_differences_2010}). Specifically, studies consistently find that women produce significantly fronter /s/ than men, and that gay men produce significantly fronter /s/ than straight men (eg. \cite{podesva_sexuality_2015, hazenberg_language_2012, flipsen_acoustic_1999}; also reviewed in \cite{calder_gay_2020}), and that /s/ is linked to women, gayness, and femininity by listeners (eg. \cite{campbell-kibler_intersecting_2011, maegaard_s_2015}).




\section{Style Shifting}
  Changes in a person's speech which arise across differences in conversational context -- eg. interlocutors, listeners, overhearers, topic, formality, and attitude -- are referred to by the term style-shifting in the field of sociolinguistics, and has been a topic of much research since the beginning (eg. NEEDCITE). Modern sociolinguistic research understands 
  
  
\textcite{kiesling_style_2009}, in attempt to ``connect the everyday use of language variation in discourse to the ways that it patterns on larger social scales'', argues ``that stancetaking is where indexicality in variation begins... any choice of linguistic form made by a sppeaaker is based ultimately on the interpersonal or epistemic stance they wish to ake with their various interlocutors at a particular time, and that it is stances that become associated, through cultural models, with various identities.'' He defines stance as ``a person's expression of their relationship to their talk (...), and a person''s expression of their relationship to their interlocutors (...).''

``Stance is at the center of this process of indexicalization because it is not each individual liguistic and social practice that an interactant decides on, but what stance to take in a particular situation.'' (179)


  
  
  Trans women, transfems, and AMAB nonbinaries styleshifting studies:
	
  
  \section{Recognition}
    \section{Queer Linguistics, Trans Linguistics}
    \section{The Variable: /s/}
	\section{Indexicality}
	\section{Intersectionality, Trans linguistics, and Raciolinguistics}
   	\section{Performativity}
   	\section{Feminism, Queer Theory, Literary theory, and Trans Studies}





\section{disjointed stuff}
	
	 Interestingly, perceptual results from cisgender heterosexual white listeners do not match results on transgender speakers; \textcite{jessee_}

	the studies do not open the possibility of investigating transgender and non-white indexicalities of /s/ on its own terms, and instead only permit analyses which echo the marginalizing judgements of the cisgender listening subject.
 

However, if these perceptual results were to be applied to the transgender speakers in \textcite{podesva_sexuality_2015, hazenberg_language_2012}, transgender speakers could only be figured as indexing their gender  insufficiently. In this way, the studies do not open the possibility of investigating transgender and non-white indexicalities of /s/ on its own terms, and instead only permit analyses which echo the marginalizing judgements of the cisgender listening subject. 

	 
	 
	 
	 One approach to this problem is the notion of indexical orders developed by  Silverstein \parencite*{silverstein_language_2023, silverstein_indexical_2003}. As linguistic forms are construed as indexical signs and gain meanings within a particular indexical field, it is possible that new meanings will arise which build off of and are informed by the original meaning. This can take place both within a single indexical field, and cross-cutting the overlaps of multiple. It is helpful here to 

For example, a particular linguistic form may come to index femininity in the hegemonic indexical field. This semiotic development would already exist informed by an innumerably long chain of previous semiotic knowledge and material, and for this reason is said to exist at the $n$th indexical order (\cite{silverstein_indexical_2003}). However, when that form is used by those percieved as men, the combination of these two significations brings the interpretation to a new indexical order -- not only is femininity indexed, but gay masculinity often arises as an additional meaning at the $n + 1$st order. This development of the gay masculinity index at the $n + 1$st order may also occupy the hegemonic indexical field, especially if users of this form are stigmatized if they are percieved as men. Yet, further indexical orders can develop that are situated within other indexical fields. For example, if one group of drag queens uses this form differently than another, an indexicality at the $n + 1 + 1$nd order may emerge whereby those in that drag community may understand the form to index membership in a particular drag group.
	 
	 
	 


	 
	 \textcite{calder_whose_2021}
	    


By making claims about the indexical meaning of linguistic forms, scholars are conjecturing based on their socioculturally mediated knowledge. 
	
	
	Although these scholars do not make value judgements of these speakers, descriptions nonetheless cannot be neutral (see \cite{zimman_beyond_2021})
	

	
	
	The eventual connection of sociolinguistic and linguistic anthropological theories of this social semiotics  to the broader social theory of performativity only came about as a result of significant development in these theories. \textcite{eckert_three_2012} breaks down this developments into three stages, or `waves', which. It is helpful to explain these stages with reference to a simple example of the meaning-making process by which the sound of a fire alarm comes to mean fire. In the first wave, scholarship noticed the existence of sociolinguistic variation, and concluded (and indeed confirmed, in some cases) that a listener might be able to notice this variation and make meaning of it. The first wave noticed that fire alarms often go off when there is a fire, and posited that listeners might understand this causal relationship and conclude that the meaning of the sound of a fire alarm is fire. The second wave acknowledged the importance of cultural differences, noting that different cultures may have fire alarms with differing sounds. In this way, the relationship be
	
	
	 had been came to integrate performativity into the theory of indexicality, which had been primarily developed in paralell by  Michael Silverst


, especiawho seek to understand how language use shapes and is shaped by society and culture, came to take up this theory, seei



Scholars of language and gender have come to understand the sounds of the voice as resources for gender performativity (\cite{bucholtz_feminist_2014}). Increasingly, research has looked to trans voices and perspectives for insight on this process, investigating the role of the voice in presentations of gender identities beyond the binary (see \cite{zimman_transgender_2020}). This research typically involves the interpretation of connections between speakers' identities and the sounds of their voices. These connections are only interpretable within a set of culturally-available associations (see \cite{calder_whose_2021}) that a listener might apply in perception of a speaker's voice. Existing studies often choose to focus on only one set of associations, and thus only one type of theoretical listener-- at present, either the set of associations of the speakers' own local queer community (eg. \cite{calder_sissy_2019}), or the hegemonic set of normative associations as represented by findings among the cis straight white population (eg. \cite{hazenberg_language_2012, podesva_sexuality_2015}). The latter type of study, by interpreting trans voices in their deviance from or compliance with cisnormativity, runs the risk of tacitly enforcing these standards, marginalizing anti-normative trans voice practices (\cite{calder_interrogating_2025}). Furthermore, focus within a single type of theoretical listener -- hegemonic or otherwise -- often sidelines the intersubjective nature of identity, by which identity emerges in the establishment and contrast of self and other, speaker and listener (see \cite{bucholtz_locating_2010}). Recent research has taken steps to place more even focus on both sides of the speaker-listener pair, investigating how trans voices are percieved by listeners (\cite{jessee_cisgender_2025}) and trans people's understandings of these perceptions (\cite{heffernan_constructing_2023}). However, research which treats the listener as a variable, investigating how trans speakers' understandings of how they will be percieved manifest in their vocal practice, has not yet been conducted. By filling this gap, I hope to demonstrate the diverse strategies trans speakers employ to resist transphobic perceptions.



\section{Indexicality and /s/}
Sociolinguistics studies the relationship between language and the socio-cultural world. This requires moving beyond the scope of traditional linguistic analysis, analyzing language in use as an interactional medium rather than language figured as an abstract system, asking not only ``what did they say'' but ``how did they say it'' and ``what were the social effects of what they said and how they said it?'' (\cite{silverstein_language_2023}). In quantitative linguistics, also termed variationist linguistics, the unit of this analysis is the sociolinguistic variable, often understood by researchers [do I need to clarify who? or is the Eckert citation enough] to be an isolated aspect of language use, at any level from phonetics to pragmatics, for which multiple variant forms fulfill similar abstract functions and not effect changes in denotational (literal) meaning, but instead effect changes in what has been termed interactional (social) meaning (\cite{eckert_three_2012}). Variants are considered to be signifiers of social meanings. 



The theory of semiotics (the study of signs) developed by \textcite{gal_signs_2019} can, however, be used to complicate and problematize this typical understanding. The definition above hides the culturally- and ideologically-mediated judgements which are involved in positing a sociolinguistic variable. 

 based upon researcher's conjecturing of a variant's indexical meaning through observed correlations

The study of language and gender has developed sophisticated theories about the discursive construction of gender through the elaboration of theories of the performative and intersubjective process of indexicality. Recent models acknowledge that linguistic variants do not have meaning in themself, but rather are construed to have meaning through a social process; listeners guess, or conjecture, that a variant holds a particular social meaning based on their ideologically-mediated knowledge of that variant's co-occurence with other signs (\cite{gal_signs_2019}).  

Unlike in traditional models, where a certain variant is claimed to have a certain meaning based upon researcher's conjecturing of a variant's indexical meaning through observed correlations, this model acknowledges that meaning arises as the result of a highly relative, contingent, and contested process (\cite{ckert_three_2012}).



Variationist production research which finds correlations between gender groups and linguistic forms conjecture that these correlations point to a system of meanings in which certain forms are conjectured as signs of gender through their percieved co-occurence. (\cite{gal_signs_2019, calder_interrogating_2025})


For \textcite{gal_signs_2019}, signs and their objects are constituted by conjectures; social actors conjecture (guess) that a phenomenon is a sign, and that it signals a particular object (meaning), using their knowledge. If I walk outside and see that the streets are wet, I might conjecture that it rained. This conjecture would rely on my knowledge, or belief, that wet streets are a sign, and that they signal rain.

Using the framework of indexicality, by which linguistic features come to index (`point to') social meanings, scholars have considered the social meaning of fronted /s/ to be femininity. That /s/ comes to be associated with gay masculinity is understood by \textcite{calder_sissy_2019} to be an example of a meaning of a higher indexical order -- a notion developed by \textcite{silverstein_indexical_2003} by which a linguistic feature comes to have new meanings through the process of adaptation to a new usage context. In this case, the use of fronted /s/ by gay men adapts its original meaning, femininity, to its new context, gayness. This structure of meanings comes to be through the imagining of gay men as more feminine, or more woman-like, than straight men. \textcite{calder_sissy_2019} sees this as an example of the ideological process \textcite{irvine_language_2000} term fractal recursivity by which a divison of society gets projected inside of itself-- gay men become ``men who are women'' through fractal recursivity, and thus, at their use of a `woman's' linguistic feature, new indexical order of /s/ emerges, going from a link from fronter /s/ to femininity to a link from fronter /s/ in men to gayness.

Although this framework of meanings works very well to describe the data on white, middle-class, cisgender westerners,  \textcite{calder_whose_2021}  points out that this structure may not always be active for those outside of this hegemonic group. If fronted /s/ were simply a feminine and gay-male index, trans people would be expected to pattern just the same as cis people. The fact that they do not (as seen in, eg. \cite{hazenberg_language_2012, podesva_sexuality_2015, zimman_hegemonic_2013}) also cannot be explained as simply another level of fractal recursivity; unlike sexuality identity, which can construe itself ideologically as a layer over and above gender and sex in the fractal recursive tree, this is untenable for transness, which is unavoidably present at all levels. Similarly, that these patterns do not hold outside of the sphere of whiteness (see eg. \cite{calder_whose_2022, jessee_cisgender_2025, maegaard_s_2015}) exposes the above understanding of the meaning structure of fronted /s/ as limited and incomplete. 

Using the concept of the indexical field, developed by \textcite{eckert_variation_2008}, which refers to the set of all meanings a linguistic feature can have, depending on context,  \textcite{calder_whose_2021} argues that the majority of the research on /s/ which does not fall into the more recent fields of raciolinguistics and trans linguistics, tacitly normalizes the white, cis, western, middle-class subset of the much more expansive indexical field. They argue that an account of how sex, gender, sexuality, (and race) are produced through language is incomplete without investigations of ``speaking and listening subjects outside of White and cisgender populations.'' (43) This further motivates my study of transfeminine speakers; I hope to find how they navigate this hegemonic meaning structure, and what meanings they use /s/ frontness to create. A more recent theorization of indexicality by \textcite{barrett_sexuality_2024} suggests that speakers and listeners bring their own indexical inventories to an interaction, and differences between these, or intentional exertion of power over another, can result in instances of misrecognition, by which the intended meaning of a n indexical behavior is misinterpreted. By looking at how transfeminine speakers navigate the possibility of misrecognition (eg. a fronted /s/ making them sound like a gay man, or a backed /s/ making them sound like a man rather than a lesbian), I hope evaluate this theoretical position.

Furthermore, I hope to raise questions about the status of `meaning' as the function of linguistic behavior. In the same way as literary theorists and feminists have attempted to move analysis in their fields away from the search for a meaning and towards an understanding of doing (\cite{deleuze_november_1987, butler_gender_1990}; see also \cite{livia_its_1997}), I feel it is possible that sociolinguistics needs to take this turn; my study of how transfeminine speakers deal with those who hope to impose meaning on their speech is an attempt to address this question: do they engage in the game of meaning-making, either by developing their own meanings, taking up up hegemonic meanings, or strategizing around the indexical inventory of their interlocutor, or do they reject `meaning' entirely, instead focusing on what their speech can do for them within whichever context they find themselves in. Are these theorizations even properly distinct?


Persuant to both the emperical gaps around transfeminine speech, and the theoretical questions just raised, I plan to investigate the \textbf{research question}: How do transfeminine speakers' pronunciations of English /s/ compare in speech targeted to a cisgender interlocutor versus in speech targeted to a transgender interlocutor with respect to frontness, and how does this vary across speakers of differing social backgrounds and identities?
       
\subsection{Indexical Inversion}
\textcite{inoue_what_2004} points out that indexing is a temporal process; the index follows, or outlasts, that which is indexed.

\textcite{jessee_cisgender_2025} use the framework of indexical inversion to understand their findings; that cisgender listeners expect hypernormativity from trans voices cannot be seen as a result of bottom-up indexical relations, by which expectations of a social group are construed from experience with the real behavior of that group-- instead, indexical relations are inverted, and expectations are based on cisnormative ideologies which seek to other and stigmatize trans people. 


    
% Questions this lit review should answer:
% - why focus on trans people?
% - are trans people's voices percieved differently than cis people's?
% - do trans people know this?
% - why focus on transfeminine people?
% - why look at /s/?
% - how do trans people behave with respect to /s/?
% - why focus on production?

% Big questions this thesis seeks to contribute to:
%- what is gender
% - what is sexuality
% - how and why do we change our speech based on who we're talking to
% - how does gender come about through everyday interaction
% - agency question
% - perception vs. production























    The long-standing commitment of linguistic research has been to a philosophy of descriptivism, often contrasted against prescriptivism (cite). This philosophy broadly aligns linguistics with other scientific disciplines, positioning itself as no different in goal and method than a natural science such as biology; just as biology aims to describe the living world, linguistics aims to describe the linguistic world. Although it perhaps sees itself as nothing but a commitment to factuality, \textcite{zimman_beyond_2021} discusses how linguistic descriptivism can be more accurately described as an apoliticizing language ideology which conveniently ignores language's deep politicality. In this view it is perhaps unsurprising that sociolinguistics, a subfield dedicated to understanding language in use (cite), has had a tendency to accept ideological constructions such as sex, gender, race and sexuality as unproblematic categories (see \cite{livia_its_1997}). For instance, this apoliticizing approach finds it much easier to ask `how does the speech of men and women differ?' than to ask the more political and critical question `how is gender discursively constructed through everyday speech?'. However, this tide has begun to turn;  three waves of indexicality eckert.
